\begin{abstract}
Vision-language models (VLMs) excel at many tasks, yet continue to struggle with spatial reasoning—problems where the key information is not directly observable in the input. Many spatial questions require \emph{imaginative perception}: simulating an unseen viewpoint, tracing a trajectory through an occluded space, or integrating partial views into a coherent spatial map. Humans naturally support this kind of reasoning through imagination. Prior work has introduced intermediate visual representations (e.g., visual thoughts, depth, or box tokens), but these intermediates often refine structure already visible rather than predicting the missing spatial structure implied by the evidence.
We introduce \textbf{Imaginative Perception Tokens (IPT)}, intermediate perceptual representations that externalize what a VLM would perceive under an alternative spatial configuration while remaining consistent with the observed input. To study this capability, we formulate three tasks that require imaginative perception: \textbf{Perspective Taking (PET)}, \textbf{Path Tracing (PT)}, and \textbf{Multiview Counting (MVC)}. For each task, we construct datasets of $\sim$20K examples spanning simulated and real-world settings, paired with ground-truth intermediate imaginations, final answers, and curated evaluation benchmarks.
Using the unified VLM BAGEL~\cite{deng2025bagel} as our backbone, IPT supervision improves spatial reasoning across several settings and often outperforms textual chain-of-thought training, even when no image is generated at inference time. For example, on MVC, IPT improves accuracy by 3.4\% and achieves performance competitive with strong closed-source models on Path Tracing. We also find that mixed training with IPT and label-only data can further improve performance. In contrast, textual chain-of-thought can be detrimental on these tasks, substantially degrading performance in some cases, highlighting a modality mismatch when forcing spatial computation through language. Overall, IPT provides a principled supervision signal for reasoning over unobserved structure, yielding stronger spatial generalization and a more interpretable intermediate aligned with the underlying geometry of the task. Code will be released at the \href{https://mahtabbigverdi.github.io/Imaginative-tokens.github.io/}{project page}.
\end{abstract}